\documentclass[a4paper,12pt]{article}
\usepackage[T2A]{fontenc}
\usepackage[utf8]{inputenc}
\usepackage[russian]{babel}
\usepackage{amsmath,amssymb,mathtools}
\renewcommand{\familydefault}{\sfdefault}
\usepackage{icomma}
\usepackage{geometry}
\geometry{left=21mm,right=21mm,top=21mm,bottom=23mm,headheight=15pt}
\usepackage{microtype}
\usepackage{xcolor}
\usepackage{tikz}
\usetikzlibrary{arrows.meta,positioning,shapes.geometric,calc}
\usepackage{tkz-euclide}
\usepackage{tabularx,booktabs}
\usepackage{enumitem}
\usepackage{fancyhdr}
\usepackage{parskip}
\usepackage{setspace}
\usepackage{hyperref}
\hypersetup{hidelinks,pdftitle={Чек-лист по математике. 12 ноября 2026},pdfauthor={Артём Александрович Лёвин}}
\definecolor{accent}{HTML}{008080}
\definecolor{darkaccent}{HTML}{005D62}
\definecolor{bodytext}{HTML}{252A2B}
\definecolor{softgray}{HTML}{667176}
\color{bodytext}
\setstretch{1.12}
\setlength{\parskip}{6pt}
\setlist[itemize]{leftmargin=18pt,itemsep=2pt,topsep=4pt}
\setlist[enumerate]{leftmargin=21pt,itemsep=6pt,topsep=4pt}
\setlength{\emergencystretch}{2em}
\pagestyle{fancy}
\fancyhf{}
\fancyhead[L]{\footnotesize\color{softgray}Практикум по математике}
\fancyhead[R]{\footnotesize\color{softgray}Ярослав Верещагин}
\fancyfoot[C]{\footnotesize\thepage}
\renewcommand{\headrulewidth}{0.3pt}
\renewcommand{\footrulewidth}{0pt}
\renewcommand{\sectionmark}[1]{}
\makeatletter
\renewcommand\section{\@startsection{section}{1}{0pt}{2.0ex plus .8ex minus .2ex}{0.9ex plus .3ex}{\normalfont\Large\bfseries\color{darkaccent}}}
\renewcommand\subsection{\@startsection{subsection}{2}{0pt}{1.6ex plus .5ex minus .2ex}{0.5ex plus .2ex}{\normalfont\normalsize\bfseries\color{darkaccent}}}
\makeatother
\newcommand{\idea}[1]{\par\noindent\textbf{\color{darkaccent}Главная мысль.} #1\par}
\newcommand{\beware}[1]{\par\noindent\textbf{\color{darkaccent}Проверьте.} #1\par}
\newcommand{\ex}[1]{\par\noindent\textbf{Пример.} #1\par}
\begin{document}

% Титульный лист
\thispagestyle{empty}
\begin{center}
\vspace*{22mm}
{\large\color{darkaccent} МАТЕМАТИКА}\par
\vspace{22mm}
{\LARGE Чек-лист}\par
\vspace{13mm}
{\LARGE\color{accent}\bfseries Вычисления, алгебра\par
и геометрия}\par
\vspace{16mm}
{\normalsize Практика задач первой части экзамена}\par
\vspace{43mm}
{\large Ярослав Верещагин}\par
\vspace{15mm}
{\normalsize 12 ноября 2026 г.}\par
\vfill
{\small Лёвин Артём Александрович\par
эксклюзивно для Ярослава Верещагина}\par
\vspace{12mm}
\begin{tikzpicture}[x=1cm,y=1cm]
 \draw[accent,line width=.8pt] (0,0) .. controls (1.0,.5) and (1.9,-.45) .. (3.3,-.03)
    .. controls (4.5,.43) and (5.2,-.34) .. (6.5,0)
    .. controls (7.7,.45) and (8.4,-.35) .. (9.6,.02);
 \fill[darkaccent] (4.8,.035) circle (1.5pt);
\end{tikzpicture}
\vspace{3mm}
\end{center}
\clearpage

\section{Числа: дроби, деление, корни}
\subsection{Деление десятичных дробей}
\idea{Сначала уберите запятые из делителя: умножьте и делимое, и делитель на одну и ту же степень десяти. Затем сократите дробь.}
\[
 \frac{4{,}2}{1{,}2}=\frac{42}{12}=\frac{7}{2}=3{,}5.
\]
\ex{При делении в столбик дробь \(\tfrac{8}{11}\) даёт периодическую десятичную запись:
\[
 \frac{8}{11}=0{,}727272\ldots=0{,}(72),\qquad 0{,}72<\frac{8}{11}<0{,}73.
\]
Если при сравнении на числовой прямой есть отрицательные дроби, сначала учитывайте знак, затем величину дроби. Например, \(-\tfrac8{11}<-\tfrac8{25}<0\).}
\beware{У каждого вычитания в столбике должен получаться верный остаток. После \(80-77=3\) следующий шаг --- \(30-22=8\), а не \(30-28\).}

\subsection{Корень из произведения степеней}
\idea{Пользуйтесь \(\sqrt{u^2}=|u|\); если знак выражения известен, модуль можно раскрыть.}
\ex{При \(a=\sqrt2\) и \(b=3\) вычислите
\[
\begin{aligned}
\sqrt{0{,}25a^4b^6}
 &=0{,}5\sqrt{(a^2b^3)^2}
  =0{,}5a^2|b|^3\\
 &=0{,}5\cdot2\cdot27=27.
\end{aligned}
\]
В данном примере \(a^2\ge0\) и \(b>0\).}
\beware{Нельзя заменять \(\sqrt{u^2}\) на \(u\) без проверки знака. В общем случае нужен модуль.}

\section{Линейное уравнение с дробью}
\idea{Чтобы избавиться от числового знаменателя, умножайте \emph{каждое} слагаемое обеих частей уравнения на этот знаменатель.}
\ex{Решим \(4x-\tfrac{x}{4}=-6\) двумя способами.}
\textbf{Способ 1. Умножение на 4:}
\[
4\cdot4x-4\cdot\frac{x}{4}=4\cdot(-6)
\ \Longrightarrow\ 16x-x=-24
\ \Longrightarrow\ \boxed{x=-\frac85}.
\]
\textbf{Способ 2. Общий знаменатель:}
\[
\frac{16x-x}{4}=-6
\ \Longrightarrow\ \frac{15x}{4}=-6
\ \Longrightarrow\ \boxed{x=-\frac85}.
\]
\beware{Знаменатель 4 --- постоянное ненулевое число, поэтому умножение на 4 равносильно. Если знаменатель содержит \(x\), сначала указывайте ограничения.}

\section{Вероятность: благоприятные исходы}
Если все участники равновероятно могут быть выбраны, то
\[
\boxed{\displaystyle P=\frac{N_{\text{благ}}}{N_{\text{всего}}}}.
\]
\ex{На олимпиаде 250 участников; в первые две аудитории распределили по 90 человек, остальные --- в третью. Вероятность выбрать участника из третьей аудитории:
\[
N_3=250-90-90=70,\qquad
P=\frac{70}{250}=\frac7{25}=0{,}28.
\]}
\beware{Не путайте \emph{число благоприятных участников} с \emph{общим числом участников}. Вероятность всегда находится между 0 и 1.}

\section{Квадратные неравенства: метод интервалов}
\idea{Найдите нули выражения, определите знак между ними и выберите интервалы, соответствующие знаку неравенства.}
\ex{Решим \(x^2-4x\le0\).}
\[
 x^2-4x=x(x-4),\qquad x_1=0,\quad x_2=4.
\]
\begin{center}
\begin{tikzpicture}[x=1.12cm,y=1cm,>=Latex]
 \draw[->,line width=.65pt] (-3.1,0)--(5.8,0) node[right]{$x$};
 \draw[line width=4pt,accent!35] (0,0)--(4,0);
 \fill[accent] (0,0) circle (2.6pt) node[below=5pt] {$0$};
 \fill[accent] (4,0) circle (2.6pt) node[below=5pt] {$4$};
 \node[above] at (-1.5,0) {$+$};
 \node[above] at (2,0) {$-$};
 \node[above] at (4.9,0) {$+$};
\end{tikzpicture}
\end{center}
Проверка знака на среднем интервале: при \(x=1\) получаем \(1-4=-3<0\). Значит, искомый участок --- от 0 до 4, \textbf{включая концы}:
\[
 \boxed{x\in[0;4]}.
\]
\beware{Совпадение нулей с рисунком ещё не гарантирует правильный вариант ответа. Проверьте знаки; у нестрогого неравенства включайте допустимые нули.}

\clearpage
\section*{Алгоритм: метод интервалов}
\thispagestyle{plain}
\vspace{5mm}
\begin{center}
\begin{tikzpicture}[
  >=Latex,
  start/.style={ellipse,draw=darkaccent,very thick,align=center,text width=7.8cm,minimum height=1.1cm,inner sep=5pt},
  act/.style={rounded corners=3pt,rectangle,draw=accent,thick,align=center,text width=9.3cm,minimum height=1.3cm,inner sep=6pt},
  decision/.style={diamond,aspect=2.9,draw=darkaccent,thick,align=center,text width=6cm,inner sep=4pt},
  branch/.style={rounded corners=3pt,rectangle,draw=accent,thick,align=center,text width=4.5cm,minimum height=1.25cm,inner sep=4pt},
  arr/.style={-{Latex[length=3mm]},thick,draw=darkaccent},
  node distance=10mm]
\node[start] (start) {Неравенство \(f(x)>0\), \(f(x)<0\), \(f(x)\ge0\) или \(f(x)\le0\)};
\node[act,below=of start] (zero) {Найдите нули выражения: решите \(f(x)=0\)};
\node[act,below=of zero] (axis) {Отметьте нули на числовой прямой, разбейте её на интервалы};
\node[act,below=of axis] (sign) {Определите знак на каждом интервале по пробной точке (или множителям)};
\node[act,below=of sign] (pick) {Выберите интервалы с нужным знаком: \(+\) для \(>0\), \(-\) для \(<0\)};
\node[decision,below=10mm of pick] (strict) {Неравенство нестрогое?};
\node[branch,below left=16mm and 11mm of strict] (yes) {Да: добавьте нули,\ если они допустимы};
\node[branch,below right=16mm and 11mm of strict] (no) {Нет: нули\ не включайте};
\node[start,below=25mm of strict] (end) {Запишите ответ объединением промежутков};
\draw[arr] (start)--(zero);
\draw[arr] (zero)--(axis);
\draw[arr] (axis)--(sign);
\draw[arr] (sign)--(pick);
\draw[arr] (pick)--(strict);
\draw[arr] (strict.south west)--node[left,pos=.45]{\small да}(yes.north);
\draw[arr] (strict.south east)--node[right,pos=.45]{\small нет}(no.north);
\draw[arr] (yes.south) |- ($(end.west)+(-.3,0)$) -- (end.west);
\draw[arr] (no.south) |- ($(end.east)+(.3,0)$) -- (end.east);
\end{tikzpicture}
\end{center}
\clearpage

\section{Движение: сумма расстояний за секунды}
\idea{Если каждую следующую секунду тело проходит на одно и то же расстояние больше, составьте последовательность и сложите \emph{все} члены за нужное число секунд.}
\ex{За первую секунду тело прошло \(0{,}6\) м, далее за каждую секунду --- на \(0{,}3\) м больше. Всего 12 секунд.}
\[
\begin{array}{c|cccccc}
\text{Секунда}&1&2&3&4&5&6\\\hline
\text{Путь, м}&0{,}6&0{,}9&1{,}2&1{,}5&1{,}8&2{,}1\\
\end{array}
\]
\[
\begin{array}{c|cccccc}
\text{Секунда}&7&8&9&10&11&12\\\hline
\text{Путь, м}&2{,}4&2{,}7&3{,}0&3{,}3&3{,}6&3{,}9
\end{array}
\]
Складывайте с краёв: \(0{,}6+3{,}9=4{,}5\), \(0{,}9+3{,}6=4{,}5\) и т.~д. Всего шесть таких пар:
\[
\boxed{S=6\cdot4{,}5=27\ \text{м}}.
\]
\beware{В таблице указан путь за \emph{отдельную} секунду, а нужен \emph{общий} путь.}

\section{Геометрия: узнаём общий элемент}
\subsection{Площади треугольников с общей высотой}
У треугольников \(ABD\) и \(ABC\), если \(D\in AC\), одна высота из вершины \(B\). Поэтому
\[
 \boxed{\displaystyle\frac{S_{ABD}}{S_{ABC}}=\frac{AD}{AC}}.
\]
\begin{center}
\begin{tikzpicture}[scale=.79]
 \tkzDefPoints{0/0/A,8/0/C,2/4/B,1.778/0/D}
 \tkzDrawPolygon[thick](A,B,C)
 \tkzDrawSegment[line width=1.0pt,color=accent](B,D)
 \tkzLabelPoints[below](A,D,C)
 \tkzLabelPoints[above](B)
 \tkzDrawPoints(A,B,C,D)
 \node[below=12pt] at (.8,0) {$8$};
 \node[below=12pt] at (4.8,0) {$28$};
\end{tikzpicture}
\end{center}
\ex{Даны \(AD=8\), \(DC=28\), \(S_{ABC}=4050\). Тогда \(AC=36\) и
\[
 \frac{S_{ABD}}{4050}=\frac8{36}
 \ \Longrightarrow\ \boxed{S_{ABD}=900}.
\]}
\beware{Отношение площадей равно отношению оснований только при \emph{одинаковых высотах}.}

\subsection{Касательные к окружности}
Радиус в точку касания перпендикулярен касательной. Если окружность касается сторон угла \(ACB=126^\circ\) в точках \(A,B\), а \(O\) --- центр, то у четырёхугольника \(AOB C\) два прямых угла:
\[
90^\circ+\angle AOB+90^\circ+126^\circ=360^\circ.
\]
Следовательно, \(\boxed{\angle AOB=54^\circ}\).

\subsection{Ромб, трапеция, параллельные прямые}
\begin{itemize}
 \item \textbf{Ромб.} Его диагонали делят углы пополам. При остром угле \(64^\circ\) угол между стороной и \emph{меньшей} диагональю равен \(\frac{180^\circ-64^\circ}{2}=\boxed{58^\circ}\).
 \item \textbf{Трапеция.} Средняя линия \(m=\frac{a+b}{2}\), площадь \(S=mh\). При \(a=5\), \(b=8\), \(h=3\) получаем \(m=6{,}5\), \(S=19{,}5\).
 \item \textbf{Перпендикулярность.} Если две \emph{различные} прямые на плоскости перпендикулярны одной и той же прямой, то они параллельны друг другу.
\end{itemize}

\section{Уравнения высокой степени и разность квадратов}
\idea{Прежде чем раскрывать большие степени, ищите формулу \(A^2-B^2=(A-B)(A+B)\).}
\ex{Решим \(x^8=(20-x)^4\):}
\[
\begin{aligned}
x^8-(20-x)^4&=0,\\
\bigl(x^4-(20-x)^2\bigr)\bigl(x^4+(20-x)^2\bigr)&=0.
\end{aligned}
\]
Вторая скобка \(x^4+(20-x)^2\) строго положительна при любом \(x\): оба слагаемых неотрицательны и не обращаются в нуль одновременно. Значит,
\[
\begin{aligned}
x^4-(20-x)^2&=0,\\
\bigl(x^2-20+x\bigr)\bigl(x^2+20-x\bigr)&=0.
\end{aligned}
\]
\(x^2-x+20=0\) действительных корней не имеет, так как \(D=1-80=-79<0\). А
\[
x^2+x-20=(x+5)(x-4)=0.
\]
\[
\boxed{x=-5\quad\text{или}\quad x=4}.
\]
\beware{Сумма квадратов \emph{не обязательно} положительна при любых данных; она равна нулю, если оба квадрата нулевые. Здесь одновременное равенство нулю невозможно.}

\clearpage
\section*{Тренировочный блок \normalfont\normalsize (путь А --- 10 задач)}
Решайте по порядку: от вычислений к уравнениям и геометрии. В задачах с промежутками записывайте ответ в виде числового множества.
\begin{enumerate}[itemsep=9pt]
\item Вычислите: \(\displaystyle\frac{6{,}3}{1{,}8}\).
\item Найдите значение \(\sqrt{0{,}36a^4b^2}\) при \(a=\sqrt5\), \(b=2\).
\item Решите уравнение \(\displaystyle 3x-\frac{x}{3}=8\) двумя способами: умножением на 3 и приведением к общему знаменателю.
\item Из 180 участников соревнования 60 находятся в первой аудитории и 48 --- во второй, остальные --- в третьей. Найдите вероятность случайно выбрать участника из третьей аудитории.
\item Решите неравенство \(x^2-5x\le0\) методом интервалов.
\item За первую секунду тело прошло \(0{,}4\) м, а за каждую следующую --- на \(0{,}2\) м больше, чем за предыдущую. Какой общий путь оно прошло за 6 секунд?
\item В треугольнике \(ABC\) точка \(D\) лежит на стороне \(AC\), \(AD=7\), \(DC=21\). Известно, что \(S_{ABC}=560\). Найдите \(S_{ABD}\).
\item Окружность с центром \(O\) касается сторон угла \(ACB=118^\circ\) в точках \(A\) и \(B\). Найдите \(\angle AOB\).
\item Острый угол ромба равен \(70^\circ\). Найдите угол между его стороной и меньшей диагональю.
\item Решите уравнение \(x^8=(12-x)^4\) разложением на множители.
\end{enumerate}

\clearpage
\section*{Ответы к тренировочному блоку}
\renewcommand{\arraystretch}{1.48}
\begin{center}
\begin{tabularx}{.94\linewidth}{@{}cX@{}}
\toprule
\textbf{№} & \textbf{Ответ} \\
\midrule
1 & \(3{,}5\) \\
2 & \(6\) \\
3 & \(x=3\) \\
4 & \(0{,}4=\frac25\) \\
5 & \(x\in[0;5]\) \\
6 & \(5{,}4\) м \\
7 & \(140\) \\
8 & \(62^\circ\) \\
9 & \(55^\circ\) \\
10 & \(x=-4\) или \(x=3\) \\
\bottomrule
\end{tabularx}
\end{center}
\subsection*{Итоговая памятка}
\begin{itemize}[itemsep=0pt,topsep=1pt]
 \item Деление с десятичными дробями: умножить числитель и знаменатель на одно и то же число, затем сократить.
 \item Корень из квадрата: проверить, нужен ли модуль.
 \item Уравнение с дробным коэффициентом: умножить \emph{все} члены на общий знаменатель.
 \item Вероятность: благоприятные исходы разделить на общее число равновозможных исходов.
 \item Неравенство: нули $\to$ знаки $\to$ нужные интервалы $\to$ включение границ.
 \item Путь за несколько секунд: сложить значения \emph{за каждую} секунду.
 \item Геометрия: искать общую высоту, перпендикулярный радиус, биссектрису диагонали или среднюю линию.
 \item Высокие степени: проверить, не скрыта ли разность квадратов.
\end{itemize}

\vspace{3mm}
\textbf{Как проверять себя:} в вычислениях возвращайтесь к исходному выражению; найденные корни подставляйте в исходное уравнение; для неравенств проверяйте точку на каждом отмеченном промежутке. В геометрии не забывайте указывать градусы и единицы площади.
\vfill
\begin{center}
{\footnotesize\color{softgray}Лёвин Артём Александрович --- эксклюзивно для Ярослава Верещагина}
\end{center}
\end{document}
